\documentclass[11pt,a4paper]{article}
\usepackage[T1]{fontenc}
\usepackage{lmodern}
\usepackage[margin=28mm]{geometry}
\usepackage{amsmath,amssymb,amsthm}
\usepackage[colorlinks=true,urlcolor=blue]{hyperref}
\newtheorem{proposition}{Proposition}
\newcommand{\Cat}{\operatorname{Cat}}
\title{A type $A_2$ counterexample to conjecture 00000000564}
\author{asfigh665}
\date{October 10, 2026}
\begin{document}
\maketitle
\begin{abstract}
Conjecture 00000000564 asserts that the specialization $F_i(1)$ of an
$F$-polynomial divides the corresponding generalized Catalan number.
In the principal-coefficient cluster algebra of type $A_2$, a single
mutation produces $F_i=1+y_1$. Its value at $y_1=y_2=1$ is $2$,
whereas the type $A_2$ generalized Catalan number is $5$.
Thus the asserted divisibility already fails in rank two.
\end{abstract}

\section{Conventions and the precise instance}
For a finite Coxeter group $W$ of rank $n$, Coxeter number $h$, and
exponents $e_1,\ldots,e_n$, the generalized Catalan number is
\[
  \Cat(W)=\prod_{j=1}^{n}\frac{h+e_j+1}{e_j+1}.
\]
Here ``generalized'' refers to the extension to finite Coxeter types;
there is no additional Fuss parameter in the conjecture.
In type $A_2$, $W\cong S_3$, $h=3$, and $(e_1,e_2)=(1,2)$.
These are equivalently the standard type $A_2$ degrees $(2,3)$.
For a direct check of the Coxeter data, the simple reflections in the
simple-root basis and their product are
\[
 s_1=\begin{pmatrix}-1&1\\0&1\end{pmatrix},\qquad
 s_2=\begin{pmatrix}1&0\\1&-1\end{pmatrix},\qquad
 C=s_1s_2=\begin{pmatrix}0&-1\\1&-1\end{pmatrix}.
\]
The matrix $C$ has order $3$ and characteristic polynomial $t^2+t+1$,
whose roots are the two nontrivial cube roots of unity. This gives
$h=3$ and exponents $1,2$; these matrix and coefficient computations
are also checked in Lean.
Consequently
\begin{equation}\label{eq:catalan}
  \Cat(A_2)=\frac{3+1+1}{1+1}\frac{3+2+1}{2+1}
           =\frac52\cdot\frac63=5.
\end{equation}

We use the standard definition of $F$-polynomials for a cluster algebra
with principal coefficients: for a cluster variable $x_{i;t}$, its
$F$-polynomial is obtained by setting all initial cluster variables
$x_1,\ldots,x_n$ equal to $1$, while retaining the coefficient variables
$y_1,\ldots,y_n$. Thus $F_i(1)$ means specialization of all coefficient
variables to $1$.

\section{One mutation gives the counterexample}
Consider the rank-two initial exchange matrix and its principal extension
\[
 B=\begin{pmatrix}0&1\\-1&0\end{pmatrix},
 \qquad
 \widetilde B=
 \begin{pmatrix}0&1\\-1&0\\1&0\\0&1\end{pmatrix}.
\]
Its Cartan companion is
$\left(\begin{smallmatrix}2&-1\\-1&2\end{smallmatrix}\right)$,
the Cartan matrix of type $A_2$. The frozen variables associated with
the last two rows are $y_1,y_2$.
For an initial seed with principal coefficients, the exchange relation
at index $k$ is
\begin{equation}\label{eq:exchange}
 x_k x_k'=
 y_k\prod_{j=1}^{n}x_j^{[b_{jk}]_+}
 +\prod_{j=1}^{n}x_j^{[-b_{jk}]_+},
 \qquad [a]_+=\max(a,0).
\end{equation}
This is also the normalized principal-coefficient relation because
$y_k\mathbin{\oplus}1=1$ in the initial tropical semifield.
Reading the first column of $\widetilde B$ and mutating at $k=1$ yields
\[
  x_1x_1'=y_1+x_2,\qquad x_1'=\frac{y_1+x_2}{x_1}.
\]
Specializing $x_1=x_2=1$ gives the actual $F$-polynomial of this
cluster variable:
\begin{equation}\label{eq:f}
  F_{1;\mu_1(t_0)}(y_1,y_2)=y_1+1,
  \qquad F_{1;\mu_1(t_0)}(1,1)=2.
\end{equation}
Only this single initial mutation is needed; there is no extrapolation
from numerical experiments.

\begin{proposition}
Conjecture 00000000564 is false for the standard finite-type generalized
Catalan number.
\end{proposition}
\begin{proof}
The seed above has finite type $A_2$. Equations~\eqref{eq:catalan} and
\eqref{eq:f} give an $F$-polynomial whose specialization equals $2$
and whose corresponding generalized Catalan number equals $5$.
An integer multiple of $2$ is even, while $5$ is odd; hence $2$ does
not divide $5$. This contradicts the asserted universal divisibility.
\end{proof}

\section{Formal verification}
The accompanying dependency-free Lean~4 project makes the initial
principal-coefficient mutation formula explicit, computes its
specialization for the displayed rank-two seed, and computes the
Catalan product using exact rational arithmetic. It then proves
nondivisibility and refutes the divisibility statement for this
particular type and mutation. A universal conjecture about all finite
types and their $F$-polynomials necessarily includes this instance.
The standard identification of the displayed Cartan matrix and Coxeter
data with type $A_2$ is stated explicitly above; the formalization uses
that concrete standard data, rather than a library of general cluster
algebras or finite Coxeter groups.
The submission README gives the build commands, theorem names,
and their correspondence with the equations in this report.

\section*{Definition reference}
The principal-coefficient definition of $F$-polynomials and the exchange
relation used here are standard. See S.~Fomin and A.~Zelevinsky,
\emph{Cluster algebras IV: Coefficients}, available at
\href{https://arxiv.org/abs/math/0602259}{arXiv:math/0602259}.
All computations needed for the counterexample have been included.
\end{document}
